
\examquestion{2025}{1}{Hyperfine structure, weak/strong Zeeman splitting and quadrupole rules}{2}
The specified electronic term \({}^2P_{1/2}\) has
\(L=1,S=1/2,J=1/2\), coupled to \(I=5/2\).

\subsection*{(a) Allowed \(F\) levels}
\(\mathbf F=\mathbf I+\mathbf J\) gives
\(F=|I-J|,\ldots,I+J\), therefore
\(\boxed{F=2,\ 3}\).

\subsection*{(b) Derive the magnetic hyperfine energies}
With \(H_{\rm hfs}=A\,\mathbf I\cdot\mathbf J/\hbar^2\),
and \(\mathbf F^2=\mathbf I^2+\mathbf J^2+
2\mathbf I\cdot\mathbf J\), take expectation values:
\[
\boxed{E_F=\frac A2[
F(F+1)-I(I+1)-J(J+1)]}.
\]
The convention absorbs \(\hbar^{-2}\) into the Hamiltonian
when angular momenta are written dimensionlessly.

\subsection*{(c) Find the splitting and sketch the diagram}
Substitute \(I(I+1)=35/4\), \(J(J+1)=3/4\):
\[
\begin{aligned}
E_{F=2}&=\frac A2\left(6-\frac{35}4-\frac34\right)
=\boxed{-\frac{7A}{4}},\\
E_{F=3}&=\frac A2\left(12-\frac{35}4-\frac34\right)
=\boxed{\frac{5A}{4}}.
\end{aligned}
\]
The signed difference is \(E_3-E_2=3A\);
the physical frequency separation is
\(\boxed{\Delta\nu=3|A|/h}\).
For \(A>0\), the \(F=3\) multiplet lies higher.
\[
\begin{tikzpicture}[x=1cm,y=.45cm]
\draw[-{Latex[length=2mm]}](0,0)--(0,7.5) node[above]{energy};
\draw[accent,very thick](.7,2)--(4,2) node[right,ink]{\(F=2:\ -7A/4\quad (5~\text{substates})\)};
\draw[accent,very thick](.7,5.3)--(4,5.3) node[right,ink]{\(F=3:\ +5A/4\quad (7~\text{substates})\)};
\draw[<->](3.3,2.1)--(3.3,5.2) node[midway,right]{\(3A\)};
\end{tikzpicture}
\]

\subsection*{(d.i--ii) Weak magnetic field}
For \(\mu_B|B|\ll |A|\) (up to numerical factors),
\(F,m_F\) remain approximately good quantum numbers:
\[
E_{Fm_F}\simeq E_F+\mu_Bg_Fm_FB,\quad
g_F\simeq g_J
\frac{F(F+1)+J(J+1)-I(I+1)}{2F(F+1)}.
\]
For \({}^2P_{1/2}\), \(g_J=2/3\); therefore
\[
\boxed{g_2=-\frac19,\quad g_3=+\frac19.}
\]
The \(F=2\) level becomes \(m_F=-2,-1,0,1,2\)
(5 lines) and \(F=3\) becomes \(-3,\ldots,3\)
(7 lines). Thus generically there are
\(\boxed{12}\) distinct weak-field substates.
Each multiplet fans into equally separated
lines in \(m_F\), but the two multiplets have
opposite slopes.

\subsection*{(d.iii) Strong magnetic field}
When \(\mu_Bg_JB\gg|A|\), the preferred basis is
\(\ket{Jm_J,Im_I}\), with two electronic branches
\(m_J=\pm1/2\), each containing six \(m_I\) values:
\[
E_{m_Jm_I}\simeq
\mu_Bg_Jm_JB+A\,m_Jm_I
\quad\text{(small nuclear Zeeman term neglected)}.
\]
The level pattern looks like two steep
approximately linear electronic fans,
each split by six hyperfine offsets; it is
not appropriate to extrapolate fixed-\(F\) lines
into this regime.

\subsection*{(e) Are higher nuclear multipoles allowed?}
The nuclear spin \(I=5/2\) \emph{can} support a
nuclear electric quadrupole (rank 2), but its
diagonal hyperfine contribution also requires a
rank-2 electronic tensor matrix element
\(\langle J\Vert T^{(2)}\Vert J\rangle\).
By Wigner--Eckart the \(3j\) triangle condition is
\(2\le J+J=1\), which fails for \(J=1/2\).
Thus the \(\boxed{\text{electric-quadrupole hyperfine
term vanishes for this level}}\), as do
higher ranks \(k>1\) at first order within
the fixed \(J=1/2\) manifold.
\solcheck{Total state count:
\(5+7=12=(2I+1)(2J+1)\).
The existence of a nuclear quadrupole alone does
not guarantee a nonzero electronic quadrupole interaction.}
